\subsection{Generación de los datos simulados}
Para cada ángulo de incidencia $\theta_i$ se calculó primero el ángulo de refracción teórico despejándolo de la ecuación \eqref{eq:snell} con $n = 1.49$:
\begin{equation}
\theta_t^{teo} = \arcsin\left( \frac{\sin\theta_i}{1.49} \right).
\label{eq:teorico}
\end{equation}
Después se multiplicó por un factor fijo de 1.08, para simular un error sistemático de alineación, y por un factor aleatorio cercano a 1 con una desviación de 2.5 por ciento, para simular el error de lectura. El resultado se redondeó a medio grado.

\subsection{Ajuste de una recta por el origen}
Como la ecuación \eqref{eq:recta} no tiene término independiente, la pendiente que mejor se ajusta a los puntos por mínimos cuadrados se calcula como
\begin{equation}
n = \frac{\sum \sin\theta_i \, \sin\theta_t}{\sum \sin^2\theta_t} = 1.396 .
\label{eq:pendiente}
\end{equation}

\subsection{Promedio y desviación estándar}
El promedio de los ocho índices de la tabla \ref{tab:angulos} y su desviación estándar se calcularon con
\begin{equation}
\bar{n} = \frac{1}{8}\sum_{k=1}^{8} n_k , \qquad
s = \sqrt{\frac{\sum (n_k - \bar{n})^2}{7}} ,
\label{eq:promedio}
\end{equation}
lo que da $\bar{n} = 1.399$ y $s = 0.023$.

\subsection{Ángulo crítico}
Sustituyendo el índice obtenido en la ecuación \eqref{eq:critico},
\begin{equation}
\theta_c = \arcsin\left( \frac{1}{1.396} \right) = 45.8 \text{ grados}.
\label{eq:critico_calc}
\end{equation}
